% !TeX spellcheck = ru_RU
% !TEX root = vkr.tex
% Опциональные добавления используемых пакетов. Вполне может быть, что они вам не понадобятся, но в шаблоне приведены примеры их использования.
\usepackage{tikz} % Мощный пакет для создание рисунков, однако может очень сильно замедлять компиляцию
\usetikzlibrary{decorations.pathreplacing,calc,shapes,positioning,tikzmark}

% Библиотека для TikZ, которая генерирует отдельные файлы для каждого рисунка
% Позволяет ускорить компиляцию, однако имеет свои ограничения
% Например, ломает пример выделения кода в листинге из шаблона
% \usetikzlibrary{external}
% \tikzexternalize[prefix=figures/]

\newcounter{tmkcount}

\tikzset{
    use tikzmark/.style={
            remember picture,
            overlay,
            execute at end picture={
                    \stepcounter{tmkcount}
                },
        },
    tikzmark suffix={-\thetmkcount}
}

\usepackage{booktabs} % Пакет для верстки "более книжных" таблиц, вполне годится для оформления результатов
% В шаблоне есть команда \multirowcell, которой нужен этот пакет.
\usepackage{multirow}
\usepackage{siunitx} % для таблиц с единицами измерений

% Для названий стоит использовать \textsc{}
\newcommand{\OCaml}{\textsc{OCaml}}
\newcommand{\miniKanren}{\textsc{miniKanren}}
\newcommand{\BibTeX}{\textsc{BibTeX}}
\newcommand{\vsharp}{\textsc{V$\sharp$}}
\newcommand{\fsharp}{\textsc{F$\sharp$}}
\newcommand{\csharp}{\textsc{C\#}}
\newcommand{\GitHub}{\textsc{GitHub}}
\newcommand{\SMT}{\textsc{SMT}}



\usepackage{makecell}

\newcommand{\gdrive}{\textsc{Google Drive}}
\newcommand{\bukyvedefont}{Buky Vede}
\newcommand{\flavuniversalfont}{Flavius Universal}
\newcommand{\ponomarunicodefont}{Ponomar Unicode}
\newcommand{\menaionunicodefont}{Menaion Unicode}
\newcommand{\unicode}{Unicode}
\newcommand{\win}{Windows-1251}


\newcommand{\json}{\textsc{JSON}}
\newcommand{\postgresql}{\textsc{PostgreSQL}}
\newcommand{\sql}{\textsc{SQL}}
\newcommand{\api}{\textsc{API}}

\newcommand{\overleaf}{\textsc{Overleaf}}
\newcommand{\vscode}{\textsc{VS Code}}
\newcommand{\python}{\textsc{Python}}
\newcommand{\vuejs}{\textsc{Vue.js}}
\newcommand{\django}{\textsc{Django}}
\newcommand{\pytest}{\textsc{pytest}}
\newcommand{\ruff}{\textsc{ruff}}
\newcommand{\docker}{\textsc{Docker}}
\newcommand{\makefile}{\textsc{Makefile}}
\newcommand{\make}{\textsc{GNU Make}}
\newcommand{\word}{\textsc{Microsoft Word}}
\newcommand{\librewriter}{\textsc{Libre Office Writer}}
\newcommand{\docx}{\textsc{docx}}
\newcommand{\pypdf}{\textsc{PyPDF2}}
\newcommand{\xelatex}{\textsc{XeLaTeX}}
\newcommand{\nginx}{\textsc{nginx}}
\newcommand{\sus}{\textsc{SUS}}
\newcommand{\systemusabilityscale}{\textsc{System Usability Scale}}
\newcommand{\gunicorn}{\textsc{gunicorn}}
\newcommand{\pdf}{\textsc{PDF}}
\newcommand{\tex}{\textsc{TeX}}



\newcommand{\todo}[1]{\textcolor{red}{TODOP: #1}}
\newcommand{\todov}[1]{\textcolor{orange}{TODOV: #1}}





\definecolor{cellgreen}{RGB}{198, 239, 206}
\definecolor{cellorange}{RGB}{255, 229, 153}
\definecolor{cellred}{RGB}{255, 199, 206}


\definecolor{eclipseGreen}{RGB}{63,127,95}
% \lstdefinelanguage{ocaml}{
% keywords={@type, function, fun, let, in, match, with, when, class, type,
% nonrec, object, method, of, rec, repeat, until, while, not, do, done, as, val, inherit, and,
% new, module, sig, deriving, datatype, struct, if, then, else, open, private, virtual, include, success, failure,
% lazy, assert, true, false, end},
% sensitive=true,
% commentstyle=\small\itshape\ttfamily,
% keywordstyle=\ttfamily\bfseries, %\underbar,
% identifierstyle=\ttfamily,
% basewidth={0.5em,0.5em},
% columns=fixed,
% fontadjust=true,
% literate={->}{{$\to$}}3 {===}{{$\equiv$}}1 {=/=}{{$\not\equiv$}}1 {|>}{{$\triangleright$}}3 {\\/}{{$\vee$}}2 {/\\}{{$\wedge$}}2 {>=}{{$\ge$}}1 {<=}{{$\le$}} 1,
% morecomment=[s]{(*}{*)}
% }

\makeatletter
\@ifclassloaded{beamer}{
    %%% Обязательные пакеты
    %% Beamer
    \usepackage{beamerthemesplit}
    \usetheme{SPbGU}
    \beamertemplatenavigationsymbolsempty
    \usepackage{appendixnumberbeamer}

    %% Локализация
    \usepackage{fontspec}
    \setmainfont{CMU Serif}
    \setsansfont{CMU Sans Serif}
    \setmonofont{CMU Typewriter Text}

    %\setmonofont{Fira Code}[Contextuals=Alternate,Scale=0.9]
    %\setmonofont{Inconsolata}
    \usepackage{polyglossia}
    \setmainlanguage{russian}
    \setotherlanguage{english}

    %% Графика
    \usepackage{pdfpages} % Позволяет вставлять многостраничные pdf документы в текст

    % Математические окружения с русским названием
    \newtheorem{rutheorem}{Теорема}
    \newtheorem{ruproof}{Доказательство}
    \newtheorem{rudefinition}{Определение}
    \newtheorem{rulemma}{Лемма}
    \usepackage{fancyvrb}
}
{}
\makeatother

\usepackage[autostyle]{csquotes} % Правильные кавычки в зависимости от языка
\usepackage{totcount}
\usepackage{setspace}
\usepackage{amsmath, amsfonts, amssymb, amsthm, mathtools} % "Адекватная" работа с математикой в LaTeX

\usepackage{fontspec}
\newfontfamily{\bukyvede}[
  Path = fonts/
]{bukyvede.ttf}

\newfontfamily{\flaviusuniversal}[
  Path = fonts/
]{FlaviusUniversal.ttf}

\newfontfamily{\ponomarunicode}[
  Path = fonts/
]{PonomarUnicode.otf}

\newfontfamily{\menaionunicode}[
  Path = fonts/
]{menaionunicode.otf}
